\section{引言}\label{sec:intro}

几何的端到端学习需要对位姿估计器求导。该估计器是一个具有多个局部极小值的非凸优化问题；一旦全局最小值的身份发生变化，其\emph{全局}解的导数就是一个不连续的对象。目前有两类方法占主导地位，而每一类都在为所得付出代价。

\textbf{确定性全局层。} SDPRLayers \cite{holmes2025sdpr} 将齐次二次约束二次规划（QCQP）提升为矩阵变量 $X\succeq0$，求解 Shor 松弛，并在松弛紧时读出一个秩一的全局解及其对偶证书。其反向传播对原始 QCQP 的 KKT 系统（而不是 SDP）求导，使用一个容忍冗余约束的隐式选择论证。该方法是可认证的，并已在矩阵加权立体定位上得到验证。其代价是结构性的：前向传播在 CPU 上求解 SDP，问题维度随变量数目增长，而保证以紧性和约束规范为条件。对称物体或近平面构型可能产生非紧松弛；此时该层没有位姿可返回。

\textbf{概率层。} EPro-PnP \cite{chen2022epropnp} 将 PnP 残差解释为 $SE(3)$ 上的玻尔兹曼分布，并用 KL 损失、以自适应多重重要性采样训练。这消除了硬分支选择：歧义位姿获得概率质量，梯度在各假设上取平均。其代价是方差：六维蒙特卡罗积分可能漏掉远处的模式，而梯度的方差必须由样本数控制。

\textbf{缺口。} 目前缺少这样一个层：（i）无需凸提升即可全局穷举；（ii）以恒定规模的反向传播可微；（iii）在分支交叉处光滑且无需采样。QPEP 这一工作线 \cite{wu2022qpep} 恰好提供了所需的代数要素：带单一单位范数约束的四元数参数化、平移的精确消去、枚举所有 KKT 根的 Gr\"obner 基求解器，以及闭式的可观测性与协方差分析。QPEP 没有提供的，是对数据的导数以及对歧义的概率处理。

本文在这些要素之上构建 \textsc{Laplace-QPEPLayer}。优化变量为 $q\in\mathbb{S}^3$；平移以闭式消去；前向传播枚举所有实 KKT 根；反向传播使用 $3\times3$ 黎曼隐式微分；歧义由枚举极小值上玻尔兹曼后验的 Laplace 近似处理。该设计刻意保持狭窄：我们不处理不等式约束（伴随线 \textsc{Constrained-QPEPLayer}）或鲁棒外点核（伴随线 \textsc{Robust-QPEPLayer}）。

\subsection{贡献}

\begin{enumerate}
\item \textbf{旋转子问题的精确代数桥梁。} 我们证明齐次旋转提升满足 $x_R(q)=2Uv_2(q)$，其中 $U\in O(10)$（定理~\ref{thm:lift}）；21 个标准 $SO(3)$ 约束恰好张成 $\nu_2(\mathbb{RP}^3)$ 的二次消失理想（定理~\ref{thm:ideal}）；并且相关的 $10\times10$ SDP 是 $\mathbb{S}^3$ 上四次型的 Veronese（Shor）松弛（定理~\ref{thm:sdp}）。适用范围是明确的：这是\emph{精确消去平移之后的旋转子问题}。它覆盖具有完整 20 维理想的 $13\times13$ SDP，我们证明后者等价于 $10\times10$ 的 SDP（定理~\ref{thm:thirteen}），但不覆盖已发布的 $24$ 约束实现——该实现的约束张成空间是 17 维，且一般而言严格更松。
\item \textbf{恒定规模的黎曼反向传播。} 我们由欧拉齐次性、黎曼 Hessian $H_{\mathbb{S}^3}=E^\top(\nabla^2J-4JI)E$ 以及 $q^\star$ 与 $t^\star$ 的隐式 Jacobian 导出 $\lambda^\star=4J$（定理~\ref{thm:implicit}），并具有精确的对映不变性：每个活跃模式一次 $3\times3$ 正定求解，而不是带冗余约束的 KKT 系统。
\item \textbf{多模态、免采样的 Laplace 层。} 平移被精确积掉；旋转通过所有枚举极小值上的 Laplace 近似处理，给出解析权重与梯度（定理~\ref{thm:laplace}）。我们给出正确的测度与归一化常数，证明远离分岔时旋转误差为 $O(\tau)$，并指出近似在何处失效。
\item \textbf{可微性证书，以及诚实的失效分析。} 我们给出 $C^\infty$ 全局位姿映射的充分条件（定理~\ref{thm:cert}），用反例说明一个自然的“精确当且仅当”命题并不成立，并隔离出需要一致渐近的分支合并机制。此外还有：一个带认证间隙的构造 Choi--Lam 实例、一个把已发布松弛与完整理想松弛分离的物理六点反例，以及与理论配套的可复现验证套件。
\end{enumerate}

\subsection{范围，以及我们不主张的内容}\label{sec:nonclaims}

本文的实证部分是对机制、正确性与偏差的受控研究，而不是感知系统的基准。具体来说：

\begin{itemize}
\item 不存在\textbf{检测器级 ADD(-S) 或 ADD-S AUC 主张}，没有基于检测器的旋转或平移误差，没有成功率，也\textbf{没有学习型前端}。第~\ref{sec:real} 节的 LINEMOD 研究直接从真实深度像素与网格顶点构造双模态代价；它是在真实几何上对歧义性、正确性与偏差/光滑性的受控测量。
\item 没有运行任何学习型或可微层基线（SDPR-IS、SDPR-SDP、EPro-PnP、BPnP）。经典 PnP 求解器（EPnP、ITERATIVE、SQPNP、IPPE）与 Theseus 仅作为同一组真实对应关系上的前向与局部参照出现；不附带任何梯度主张。Monte Carlo 基线对离散模式权重采样，并不是 EPro-PnP 的 $SE(3)$ 自适应重要性采样器。
\item UTIAS In-The-Dark 未运行。端到端训练、GPU 批量根枚举、吞吐量与峰值内存均未测量。根枚举的完备性是一个假设（A2），由残差验证而非认证。
\item 认证的 SDP 间隙是数值的（求解器容差下的 PSD 对偶矩阵），而非符号的；并且构造的 Choi--Lam 目标尚不知道能否实现为某个真实对应集的 Gram 矩阵。
\item \textbf{真实数据上的结果是取舍，而不是精度胜利。} 在 15 个真实的近平局帧上，硬 argmin 是准确的（中位数 $2.06^\circ$、$3.7$ mm ADD-S），而 Laplace 后验均值偏向翻转（中位数 $44.24^\circ$、$13.9$ mm、配对偏差 $+36.9^\circ$）。早期出现的 $58.28^\circ$ 这一数字，是把测地线公式作用在\emph{未投影}的 9 维向量平均 $\sum_kw_kR_k$ 上得到的，而后者不是旋转；它不是旋转误差，本文任何地方都不使用它。
\end{itemize}

每个数值主张都标注三种证据级别之一：\emph{论文声明}（在第~\ref{sec:formulation}--\ref{sec:method} 节中给出、且未重新测量）、\emph{独立复现（本仓库）}（由本仓库中的脚本产出，其网格、种子与容差记录在指定的结果文件中；可在此处重跑，而不是对外部论文数字的复现），或\emph{未验证}（没有测量的设计目标）。没有任何论文声明被当作独立复现来呈现。
